\chapter{Apprendre sans étiquettes : l'auto-supervision}
\label{chap:auto-supervision}

\begin{quote}\itshape
Au chapitre précédent, la machine avait toujours la bonne réponse écrite à
côté de la question. Nous allons voir maintenant que, pour les vidéos, cette
bonne réponse n'existe pas et n'existera jamais. Nous allons donc apprendre
à la machine un nouveau jeu : cacher elle-même un morceau de ce qu'elle
regarde, et essayer de le deviner. Personne ne lui souffle rien. Et
pourtant, elle apprend.
\end{quote}

\begin{prerequis}
Les chapitres~\ref{chap:image-video} et~\ref{chap:apprendre}. Il faut savoir
qu'une image est une grille de nombres et connaître sa forme, et il faut
avoir compris la boucle en cinq gestes : je prédis, je mesure mon erreur, je
calcule les pentes, je bouge les boutons, je recommence. Nous n'aurons besoin
de rien d'autre. Aucun mot nouveau ne sera employé sans être expliqué juste
avant.
\end{prerequis}

Reprenons la machine à bonbons du chapitre précédent. Pour qu'elle apprenne,
il nous fallait trois choses : les questions (le nombre de pièces), les
réponses (le nombre de bonbons), et la boucle qui corrige les boutons. Sans
les réponses, la deuxième étape de la boucle s'arrête net : impossible de
mesurer un écart, donc impossible de savoir dans quel sens tourner quoi que
ce soit.

Tout ce chapitre part de cette question, et d'elle seule : \emph{et si
personne n'écrit les réponses ?}

\section{La bonne réponse écrite à côté : l'étiquette}

\begin{enclair}
Une étiquette, c'est le petit mot qu'un humain colle sur une donnée pour
dire ce qu'elle contient. « Ceci est un chat. » Sans ce petit mot, la
méthode du chapitre précédent ne fonctionne pas.
\end{enclair}

Quand nous avons entraîné la machine à bonbons, nous connaissions les trois
réponses : deux bonbons pour Léa, cinq pour Sami, huit pour Nina. Ces trois
nombres n'étaient pas tombés du ciel. Trois enfants les avaient comptés, et
quelqu'un les avait notés sur un papier. Ce travail de notation a un nom.

\begin{definition}[L'étiquette][etiquette]
Une \emph{étiquette} est la bonne réponse que l'on écrit à côté d'une donnée,
pour dire à la machine ce qu'elle aurait dû répondre. Sur une photo, ce peut
être le mot « chat ». Sur un enregistrement de voix, ce peut être la phrase
prononcée. Sur une image médicale, ce peut être « sain » ou « malade ». Une
étiquette est toujours écrite par quelqu'un : elle ne se fabrique pas toute
seule.
\end{definition}

\begin{definition}[L'apprentissage supervisé][supervise]
On parle d'\emph{apprentissage supervisé} quand on entraîne une machine avec
des données accompagnées de leurs étiquettes. Le mot « supervisé » veut dire
« surveillé » : à chaque essai, quelqu'un — au travers de l'étiquette qu'il a
écrite à l'avance — corrige la machine. C'est exactement ce que nous avons
fait au chapitre~\ref{chap:apprendre}.
\end{definition}

\begin{exemple}[Reconnaître un chat d'un chien]
Vous voulez une machine qui, devant une photo, dise « chat » ou « chien ».
Vous rassemblez dix mille photos. Puis, une par une, vous les ouvrez, vous
regardez, et vous tapez le bon mot à côté. Dix mille fois. Quand c'est fini,
vous avez dix mille étiquettes, et vous pouvez lancer la boucle du chapitre
précédent : la machine prédit, on compare sa prédiction à l'étiquette, on
mesure l'écart, on corrige les boutons.

Retenez bien le prix payé ici : \emph{dix mille regards humains}. C'est ce
prix-là que nous allons regarder de plus près.
\end{exemple}

Tant qu'on parle de dix mille photos, tout va bien. Une petite équipe y passe
quelques semaines et le travail est fait. Voyons ce qui se passe quand on
change d'échelle.

\section{Le mur : ce que coûte vraiment une étiquette}

\begin{enclair}
Étiqueter est un travail humain. Un travail humain prend du temps. Et quand
le nombre de données devient très grand, ce temps devient si long qu'aucune
entreprise, aucun pays, aucune génération ne pourrait le fournir.
\end{enclair}

Faisons le calcul ensemble, avec des nombres simples et vérifiables.

Imaginons la tâche la plus facile qui soit : regarder une photo et taper un
seul mot. Une personne rapide met environ cinq secondes par photo, en
comptant le temps de regarder et le temps de taper. Maintenant, comptons le
temps de travail d'une personne dans une année : disons sept heures par jour,
deux cent vingt jours par an — les week-ends et les congés en moins. Cela
fait $7 \times 220 = 1\,540$ heures de travail par an.

\begin{exemple}[Étiqueter cent quarante-deux millions d'images]
V-JEPA~2.1 s'entraîne, entre autres, sur un ensemble de $142\,000\,000$
d'images. Étiquetons-les à cinq secondes pièce :

\[
  142\,000\,000 \times 5 = 710\,000\,000 \text{ secondes.}
\]

Convertissons en heures en divisant par $3\,600$ :

\[
  \frac{710\,000\,000}{3\,600} \approx 197\,000 \text{ heures.}
\]

Puis en années de travail d'une personne, en divisant par $1\,540$ :

\[
  \frac{197\,000}{1\,540} \approx 128 \text{ ans.}
\]

Cent vingt-huit ans de travail, pour une seule personne, sans jamais
s'arrêter. Et c'était la tâche la plus \emph{facile} : un mot par image.
\end{exemple}

Et les vidéos ? La plus grosse source de V-JEPA~2.1 représente à elle seule
$1\,600\,000$ heures de film. Notez bien : il ne s'agit plus du temps qu'il
faut pour étiqueter, mais du temps qu'il faudrait rien que pour
\emph{regarder} ces vidéos une fois, à vitesse normale.

\[
  \frac{1\,600\,000}{1\,540} \approx 1\,039 \text{ ans.}
\]

Plus de mille ans de visionnage. Il faudrait avoir commencé à regarder à
l'époque des premiers châteaux forts pour finir aujourd'hui — et on n'aurait
encore rien écrit.

\begin{surlesdonnees}
Le mur devient encore bien plus haut dès qu'on demande une étiquette plus
fine qu'un simple mot. Pour apprendre à une machine à dire, \emph{pour chaque
pixel}, s'il appartient à la route, au trottoir ou à un piéton, il faut
colorier l'image à la main, contour par contour. Sur des photos de rue, ce
travail prend de l'ordre d'une heure et demie par image.

Sur nos cent quarante-deux millions d'images, cela ferait
$142\,000\,000 \times 1{,}5 = 213\,000\,000$ heures, soit environ
$213\,000\,000 / 1\,540 \approx 138\,000$ ans. Cent trente-huit mille ans de
coloriage. Autant dire : jamais.
\end{surlesdonnees}

La figure~\ref{fig:montagne-heures} met ces trois durées côte à côte. Le
troisième bâton ne tient pas sur la page, et c'est bien le message du dessin.

\begin{schema}{Le temps de travail humain qu'il faudrait pour étiqueter les
données de V-JEPA~2.1, comparé à une carrière entière. Le premier bâton est
tout ce qu'une personne peut fournir dans sa vie. Le deuxième est le prix
d'un simple mot par image. Le troisième, celui du coloriage pixel par pixel,
sort de la page : il faudrait mille fois plus de temps que ce que la feuille
peut montrer.}[montagne-heures]
  \begin{scope}[x=1cm, y=1cm]
    \draw[draw=mbSlate!45, line width=0.7pt] (0,-0.25) -- (0,3.45);
    % Les trois légendes sont posées de la même façon : juste au-dessus de
    % leur bâton, alignées à gauche sur l'axe.
    % --- une carrière ---
    \node[font=\sffamily\scriptsize, text=mbSlate, anchor=west]
         at (0.05,3.15) {une carrière entière : $40$ ans};
    \draw[draw=mbTeal, fill=mbTeal!25, line width=0.8pt]
          (0,2.4) rectangle (0.9,2.9);
    % --- un mot par image ---
    \node[font=\sffamily\scriptsize, text=mbSlate, anchor=west]
         at (0.05,2.05) {étiqueter les images, un mot chacune : $128$ ans};
    \draw[draw=mbAmber!85!black, fill=mbAmber!25, line width=0.8pt]
          (0,1.3) rectangle (2.88,1.8);
    % --- le coloriage : le bâton sort de la page ---
    \node[font=\sffamily\scriptsize, text=mbSlate, anchor=west]
         at (0.05,0.95) {les colorier pixel par pixel : $138\,000$ ans};
    \draw[draw=mbRose, fill=mbRose!25, line width=0.8pt]
          (0,0.2) rectangle (7.9,0.7);
    \draw[draw=white, line width=1.6pt] (7.45,0.1) -- (7.6,0.8);
    \draw[draw=white, line width=1.6pt] (7.7,0.1) -- (7.85,0.8);
    \draw[fleche accent] (8.0,0.45) -- (8.7,0.45);
    \node[annotation, text width=2.4cm, anchor=west] at (8.85,0.45)
         {et cela continue\\ bien au-delà};
    \node[font=\sffamily\scriptsize\itshape, text=mbSlate, anchor=west]
         at (0.05,-0.65) {longueur des bâtons : temps de travail d'une seule
                          personne};
  \end{scope}
\end{schema}

Vous voyez le problème. Nous savons faire apprendre une machine, mais notre
méthode réclame un carburant — les étiquettes — que nous ne pouvons pas
produire en quantité suffisante. Or les données brutes, elles, sont partout :
des milliards de photos et de vidéos existent déjà, sans le moindre petit mot
écrit à côté.

\begin{aretenir}
Les données brutes sont abondantes et presque gratuites. Les étiquettes sont
rares et très chères. Toute la suite du livre découle de ce déséquilibre :
il faut trouver un moyen d'apprendre à partir de ce qui est abondant.
\end{aretenir}

La question devient donc : peut-on fabriquer la bonne réponse \emph{sans
personne} ? La réponse est oui, et elle est étonnamment simple.

\section{L'idée qui change tout : cacher un morceau}

\begin{intuition}
Si je cache un bout de quelque chose que j'ai sous les yeux, je connais déjà
la bonne réponse : c'est le bout que je viens de cacher. Je peux donc me
poser une question à moi-même, et me corriger tout seul. Aucun humain n'est
nécessaire.
\end{intuition}

Faisons l'expérience sur vous, tout de suite. Lisez cette phrase :

\begin{center}
\itshape « Le chat boit du \underline{\hphantom{xxxxxxx}}. »
\end{center}

Vous avez pensé « lait ». Personne ne vous a donné la réponse. Vous l'avez
trouvée parce que vous connaissez les chats, les boissons, et la façon dont
les phrases sont construites. Autrement dit : \emph{pour remplir le trou, il
a fallu comprendre le reste}. C'est toute l'idée.

Maintenant, regardez ce que cela donne si l'on cache un morceau d'une image.
Si je masque le museau d'un chien sur une photo, vous pouvez encore le
dessiner de tête, à peu près : vous savez où il va, de quelle couleur il est,
quelle forme il a. Pour cela, vous vous êtes servi des oreilles, des yeux, de
la fourrure autour. Vous avez compris la photo.

\begin{definition}[L'apprentissage auto-supervisé][ssl]
On parle d'\emph{apprentissage auto-supervisé} quand la machine fabrique
elle-même ses questions et ses réponses à partir des données brutes, sans
aucune étiquette écrite par un humain. Le préfixe « auto » veut dire
« par elle-même » : c'est elle qui se surveille. On l'abrège souvent en
\emph{SSL}, initiales de l'anglais \emph{self-supervised learning}.
\end{definition}

La figure~\ref{fig:etiquette-vs-ssl} met les deux façons de faire côte à côte.
Regardez surtout la différence de rythme entre la colonne de gauche et celle
de droite.

\begin{schema}{Deux façons d'obtenir la bonne réponse. À gauche,
l'apprentissage supervisé : un humain regarde chaque image et écrit une
étiquette, une par une. À droite, l'apprentissage auto-supervisé : la machine
recouvre elle-même un morceau de l'image, et le morceau caché devient la
bonne réponse. À gauche, il faut un humain à chaque image ; à droite,
personne.}[etiquette-vs-ssl]
  \begin{scope}[x=0.94cm, y=0.94cm]
    %================ colonne de gauche : supervisé ================
    \node[font=\sffamily\bfseries\small, text=mbSlate] at (2.2,4.95)
         {Apprentissage supervisé};
    % le personnage : une tête et un buste très simples
    \begin{scope}[shift={(0.2,2.6)}]
      \draw[draw=mbSlate, fill=mbMist, line width=0.8pt] (0,0.55) circle (0.32);
      \draw[draw=mbSlate, fill=mbMist, line width=0.8pt]
            (-0.42,-0.5) -- (0.42,-0.5) -- (0.3,0.18) -- (-0.3,0.18) -- cycle;
      \node[font=\sffamily\scriptsize, text=mbSlate] at (0,-0.95) {un humain};
    \end{scope}
    % trois images à étiqueter
    \foreach \mbk/\mby in {0/3.6, 1/2.6, 2/1.6}{
      \draw[draw=mbSlate!55, fill=white, line width=0.7pt]
            (1.6,\mby-0.35) rectangle (2.5,\mby+0.35);
      \draw[fleche] (0.75,\mby) -- (1.5,\mby);
      \node[font=\sffamily\scriptsize, text=mbInk] at (3.5,\mby)
           {« chat »};
      \draw[fleche] (2.6,\mby) -- (2.95,\mby);
    }
    \node[annotation, text width=4cm] at (2.2,0.7)
         {une image, un regard, un mot\\ puis on recommence
          142 millions de fois};
    % séparateur
    \draw[draw=mbSlate!30, dashed, line width=0.7pt] (5.2,0.3) -- (5.2,5.25);
    %================ colonne de droite : auto-supervisé ============
    \node[font=\sffamily\bfseries\small, text=mbSlate] at (8.6,4.95)
         {Apprentissage auto-supervisé};
    % l'image complète
    \draw[draw=mbSlate!55, fill=white, line width=0.7pt]
          (6.3,2.85) rectangle (7.5,3.75);
    \node[font=\sffamily\scriptsize, text=mbSlate] at (6.9,2.55)
         {l'image};
    \draw[fleche accent] (7.65,3.3) -- (8.35,3.3);
    % l'image masquée
    \draw[draw=mbSlate!55, fill=white, line width=0.7pt]
          (8.5,2.85) rectangle (9.7,3.75);
    \fill[mbSlate!35] (8.8,3.05) rectangle (9.4,3.55);
    \node[font=\sffamily\scriptsize, text=mbSlate, align=center,
          text width=3.6cm] at (9.1,4.15)
         {la machine recouvre un carré};
    % la réponse
    \draw[fleche] (9.85,3.3) -- (10.6,3.3);
    \draw[draw=mbAmber!85!black, fill=mbAmber!18, line width=0.9pt]
          (10.7,3.05) rectangle (11.3,3.55);
    \node[annotation, text width=3.2cm] at (11.1,2.25)
         {le morceau caché\\ \emph{est} la bonne réponse};
    \node[annotation, text width=4.6cm] at (8.9,1.0)
         {aucun humain n'intervient :\\ on peut le faire des milliards
          de fois, jour et nuit};
  \end{scope}
\end{schema}

Ce petit changement de point de vue a une conséquence énorme, et il faut
prendre le temps de la mesurer. Dans la colonne de gauche, le nombre
d'exemples d'entraînement est limité par le nombre d'heures que des humains
acceptent d'y passer. Dans la colonne de droite, il est limité par le nombre
de vidéos qui existent et par la puissance des machines. Ce ne sont pas du
tout les mêmes limites.

\begin{exemple}[Le même carré, deux rôles]
Prenons une image de $8 \times 8$ cases. Je recouvre les seize cases du coin
en haut à droite. Ce que je donne à la machine, c'est l'image avec son trou.
Ce que je lui demande, c'est de reconstituer les seize cases manquantes. Et
comme je viens de les recouvrir moi-même, je les connais parfaitement : je
peux donc comparer sa réponse à la vérité et calculer une erreur, exactement
comme au chapitre~\ref{chap:apprendre}.

Le même carré joue donc deux rôles à la fois : il est \emph{la question}
quand on le cache, et \emph{la réponse} quand on le compare. Voilà pourquoi
on n'a plus besoin de personne.
\end{exemple}

Ce jeu de cache-cache porte un nom, et il faut le connaître, car nous allons
l'employer jusqu'à la fin du livre.

\section{La tâche prétexte : un jeu que l'on jette après}

\begin{definition}[La tâche prétexte][pretexte]
Une \emph{tâche prétexte} est un exercice que l'on invente pour faire
travailler la machine, alors que la réponse à cet exercice ne nous intéresse
pas du tout. Le mot « prétexte » a ici son sens courant : une bonne excuse. Ce
qui nous intéresse, ce n'est pas que la machine sache remplir des trous ; c'est
tout ce qu'elle a été \emph{obligée} de comprendre pour y arriver.
\end{definition}

Cette idée surprend souvent au début, alors prenons une comparaison de la vie
de tous les jours.

\begin{intuition}
Un professeur de musique demande à son élève de jouer des gammes pendant des
mois. Personne n'a jamais acheté de billet pour aller écouter des gammes. La
gamme est un prétexte : ce que le professeur veut vraiment, ce sont des
doigts agiles et une oreille juste. Le jour du concert, l'élève ne jouera
aucune gamme, mais tout ce qu'il jouera en profitera.

Pour notre machine, c'est pareil. Remplir des trous est la gamme. Reconnaître
une action, mesurer une distance, suivre un objet : voilà le concert.
\end{intuition}

Les chercheurs ont inventé beaucoup de gammes différentes. La
figure~\ref{fig:pretextes} en montre quatre parmi les plus connues. Elles ont
toutes le même point commun : la bonne réponse se fabrique toute seule, en
partant de l'image entière.

\begin{schema}{Quatre tâches prétextes. Dans les quatre cas, on abîme
volontairement l'image d'une façon dont on garde le secret, et l'on demande à
la machine de retrouver ce qui a été fait. Comme c'est nous qui avons abîmé
l'image, nous connaissons toujours la bonne réponse, sans avoir eu besoin de
la moindre étiquette.}[pretextes]
  \begin{scope}[x=0.87cm, y=0.87cm]
    \foreach \mbx/\mbtitre/\mbquestion in {%
      0/{le puzzle}/{remets les morceaux\\ dans l'ordre},
      3.6/{la rotation}/{de quel angle\\ ai-je tourné ?},
      7.2/{la couleur}/{remets les couleurs\\ sur le gris},
      10.8/{le trou}/{devine ce qui\\ manque}}{
      \draw[draw=mbSlate!55, fill=white, line width=0.7pt]
            (\mbx,0) rectangle (\mbx+2.4,1.8);
      \node[font=\sffamily\small\bfseries, text=mbBlue] at (\mbx+1.2,2.15)
           {\mbtitre};
      \node[font=\sffamily\scriptsize, text=mbSlate, align=center,
            text width=2.6cm] at (\mbx+1.2,-0.75) {\mbquestion};
    }
    % 1. puzzle : quatre cases mélangées, numérotées dans le désordre
    \foreach \mbi/\mbj/\mbn in {0/0/3, 1/0/1, 0/1/4, 1/1/2}{
      \draw[draw=mbSlate!45, fill=mbBlue!8, line width=0.5pt]
            (0.25+\mbi*1.05, 0.2+\mbj*0.75)
            rectangle (1.2+\mbi*1.05, 0.85+\mbj*0.75);
      \node[font=\sffamily\scriptsize, text=mbSlate]
           at (0.72+\mbi*1.05, 0.52+\mbj*0.75) {\mbn};
    }
    % 2. rotation : une petite maison penchée, et la flèche qui tourne
    %    reste à l'intérieur du cadre
    \begin{scope}[shift={(4.8,0.72)}, rotate=-35, scale=0.82]
      \draw[draw=mbTeal, fill=mbTeal!12, line width=0.8pt]
            (-0.45,-0.4) rectangle (0.45,0.25);
      \draw[draw=mbTeal, fill=mbTeal!22, line width=0.8pt]
            (-0.58,0.25) -- (0,0.7) -- (0.58,0.25) -- cycle;
    \end{scope}
    \draw[fleche accent] (4.15,1.42) to[bend left=28] (5.45,1.42);
    % 3. couleur : la même petite image, grise à gauche, colorée à droite
    \fill[mbSlate!25] (7.45,0.2) rectangle (8.35,1.6);
    \fill[mbAmber!45] (8.35,0.2) rectangle (9.25,1.6);
    \foreach \mbj in {0.55,0.9,1.25}{
      \draw[draw=white, line width=0.5pt] (7.45,\mbj) -- (9.25,\mbj);
    }
    \draw[draw=white, line width=0.5pt] (7.9,0.2) -- (7.9,1.6);
    \draw[draw=white, line width=0.5pt] (8.8,0.2) -- (8.8,1.6);
    \draw[draw=mbSlate!70, line width=0.8pt] (8.35,0.2) -- (8.35,1.6);
    \draw[draw=mbSlate!45, line width=0.5pt] (7.45,0.2) rectangle (9.25,1.6);
    % 4. trou : un carré bien franc posé au milieu de l'image
    \fill[mbBlue!10] (11.05,0.2) rectangle (12.95,1.6);
    \draw[draw=mbSlate!45, line width=0.5pt] (11.05,0.2) rectangle (12.95,1.6);
    \fill[mbSlate!60] (11.7,0.55) rectangle (12.4,1.15);
    \node[font=\sffamily\small\bfseries, text=white] at (12.05,0.85) {?};
  \end{scope}
\end{schema}

V-JEPA~2.1 a choisi la quatrième : le trou. On recouvre une partie de la
vidéo et l'on demande de deviner ce qui se trouvait dessous. Nous verrons
dans la partie~II que ce choix n'est pas un hasard : c'est celui qui oblige le
plus la machine à comprendre en même temps ce qu'il y a dans la scène et
comment cela bouge.

\begin{piege}
Ne cherchez jamais à juger un modèle auto-supervisé sur sa tâche prétexte.
Une machine qui remplit très bien les trous n'est pas forcément une machine
utile ; et une machine qui les remplit imparfaitement peut être excellente
pour tout le reste. C'est même parfois le contraire de ce qu'on croit : à la
fin de ce chapitre, notre petit programme atteindra un plancher d'erreur qu'il
est \emph{impossible} de franchir, et ce sera une très bonne nouvelle.

La seule question qui compte est celle-ci : \emph{ce que la machine a appris
en jouant sert-il à autre chose ?}
\end{piege}

Ce « quelque chose » que la machine garde après avoir joué, il faut
maintenant lui donner un nom. C'est le vrai produit de tout le travail, et
c'est le sujet de la section suivante.

\section{Ce que l'on garde vraiment : la représentation}

\begin{enclair}
Quand la machine regarde une image, elle en fabrique un petit résumé en
nombres. Ce résumé s'appelle une représentation. C'est lui, et lui seul, que
l'on garde à la fin. Tout le reste — le jeu des trous, les réponses données —
part à la poubelle.
\end{enclair}

Repensez à ce que vous faites quand un ami vous raconte un film. Il ne vous
récite pas les deux heures image par image. Il vous dit : « c'est l'histoire
d'un pêcheur qui cherche son fils ». En quelques mots, il a gardé ce qui
compte et jeté le reste. Ce petit paquet d'informations essentielles, c'est
exactement ce que la machine fabrique, sauf qu'elle l'écrit avec des nombres
au lieu de mots.

\begin{definition}[L'encodeur][encodeur]
Un \emph{encodeur} est la partie du modèle qui prend une image ou une vidéo
en entrée, et qui en ressort un résumé beaucoup plus petit. On peut le voir
comme un entonnoir : beaucoup de nombres entrent, peu de nombres sortent.
C'est la pièce maîtresse de tout le livre : à la fin, c'est l'encodeur, et
presque rien d'autre, que l'on met en production.
\end{definition}

\begin{definition}[La représentation][representation]
La \emph{représentation} d'une image est la liste de nombres que l'encodeur
produit à partir d'elle. On dit aussi les \emph{caractéristiques}, ou en
anglais les \emph{features}. Deux images qui se ressemblent doivent donner
deux représentations proches ; deux images très différentes, deux
représentations éloignées.
\end{definition}

\begin{definition}[L'espace latent][latent]
L'ensemble de tous les résumés possibles s'appelle l'\emph{espace latent}.
« Latent » veut dire « caché », « qui ne se voit pas directement » : ces
nombres ne sont ni des pixels ni des mots, ce sont les notes internes de la
machine. Dire qu'une représentation « vit dans l'espace latent », c'est
simplement dire que c'est une liste de nombres inventée par le modèle, et non
quelque chose que l'on peut afficher tel quel à l'écran.
\end{definition}

La figure~\ref{fig:entonnoir} montre l'encodeur au travail.

\begin{schema}{L'encodeur, dessiné comme un entonnoir. Une image entre par le
haut, avec ses centaines de milliers de nombres. Un petit paquet de nombres
sort par le bas : c'est la représentation. Ce petit paquet vit dans ce qu'on
appelle l'espace latent. Tout l'enjeu du livre est de le rendre aussi riche
que possible.}[entonnoir]
  \begin{scope}[x=1cm, y=1cm]
    % l'image d'entrée, dessinée en petite grille
    \foreach \mbi in {0,...,5}{
      \foreach \mbj in {0,...,3}{
        \pgfmathsetmacro{\mbton}{8 + 12*mod(\mbi+2*\mbj,4)}
        \draw[draw=mbSlate!35, fill=mbBlue!\mbton, line width=0.3pt]
              (\mbi*0.42, 3.9+\mbj*0.42) rectangle
              (\mbi*0.42+0.42, 3.9+\mbj*0.42+0.42);
      }
    }
    \node[font=\sffamily\scriptsize, text=mbSlate, align=center]
         at (-1.9,4.7) {une image\\ $196\,608$ nombres};
    \draw[fleche] (-0.75,4.7) -- (-0.15,4.7);
    % l'entonnoir
    \draw[draw=mbBlue, fill=mbBlue!8, line width=1pt]
          (0,3.6) -- (2.52,3.6) -- (1.6,1.7) -- (0.92,1.7) -- cycle;
    \node[font=\sffamily\scriptsize\bfseries, text=mbBlue] at (1.26,3.25)
         {l'encodeur};
    % la sortie : une petite barre de nombres
    \foreach \mbk/\mbv in {0/0.31, 1/-0.72, 2/1.04, 3/0.15, 4/-0.48}{
      \draw[draw=mbAmber!85!black, fill=mbAmber!14, line width=0.7pt]
            (-0.35+\mbk*0.72, 0.85) rectangle (0.32+\mbk*0.72, 1.35);
      \node[font=\sffamily\scriptsize, text=mbInk]
           at (-0.015+\mbk*0.72, 1.1) {\mbv};
    }
    \draw[fleche accent] (1.26,1.62) -- (1.26,1.42);
    \node[font=\sffamily\scriptsize\bfseries, text=mbAmber!75!black,
          align=center] at (1.26,0.45) {la représentation};
    \node[annotation, text width=3.6cm] at (5.7,1.1)
         {beaucoup moins de nombres,\\ mais ceux qui comptent};
    \draw[trait annotation] (3.9,1.1) -- (3.15,1.1);
    \node[annotation, text width=3.4cm] at (5.7,4.4)
         {rien n'est étiqueté ici :\\ c'est une image brute};
    \draw[trait annotation] (4.0,4.4) -- (2.75,4.4);
  \end{scope}
\end{schema}

Mettons des nombres sur ce dessin, car c'est plus parlant.

\begin{exemple}[De deux cent mille nombres à mille]
Une image en couleur de $256 \times 256$ pixels contient
$256 \times 256 \times 3 = 196\,608$ nombres : nous l'avions calculé au
chapitre~\ref{chap:image-video}. Supposons que l'encodeur la résume en
$1\,024$ nombres. Le rapport vaut

\[
  \frac{196\,608}{1\,024} = 192 .
\]

Le résumé est donc cent quatre-vingt-douze fois plus petit que l'image. Toute
la difficulté du métier tient dans cette division : il faut jeter cent
quatre-vingt-onze nombres sur cent quatre-vingt-douze, et jeter les bons.
\end{exemple}

\begin{remarque}
Vous vous demandez peut-être : pourquoi jeter ? Pourquoi ne pas tout garder ?
Parce qu'un résumé qui garde tout n'est pas un résumé : c'est une copie, et
une copie n'a rien compris. La contrainte de place est précisément ce qui
force la machine à distinguer l'important de l'accessoire. Dans le résumé
d'un film, « un pêcheur cherche son fils » tient en six mots parce qu'on a
renoncé à la couleur des rideaux.
\end{remarque}

Nous avons donc un jeu (la tâche prétexte), et un produit (la représentation).
Il nous manque encore une chose : un moyen de savoir si ce produit est bon.

\section{Comment savoir si une représentation est bonne}

\begin{enclair}
On bloque l'encodeur pour qu'il ne bouge plus, on branche derrière lui une
toute petite machine à apprendre, et on regarde jusqu'où elle arrive. Si
même une machine minuscule s'en sort bien, c'est que le résumé était déjà
bien rangé.
\end{enclair}

Voici l'idée, en deux gestes. Premier geste : on décide que l'encodeur ne
changera plus. Ses boutons sont fixés une fois pour toutes ; on dit qu'on le
\emph{gèle}.

\begin{definition}[L'encodeur gelé][gele]
Geler un encodeur, c'est décider que ses paramètres ne seront plus modifiés.
On s'en sert alors comme d'une machine toute faite : on lui donne une image,
il rend une représentation, et rien à l'intérieur ne bouge. C'est ainsi que
V-JEPA~2.1 est évalué du début à la fin du papier — et c'est aussi, nous le
verrons dans la partie~VII, ce qui rend son déploiement, c'est-à-dire sa mise
au travail pour de vrai, si peu coûteux.
\end{definition}

\begin{exemple}[Ce que le gel permet]
Imaginons que vous ayez dix mille photos et trois questions différentes à
leur poser : à quelle distance est chaque objet, où sont les voitures, et de
quelle sorte de scène il s'agit. Comme l'encodeur ne bouge plus, vous lui
faites lire vos dix mille photos une seule fois, vous rangez les dix mille
résumés dans un fichier, et vous n'y touchez plus. Vos trois questions se
règlent ensuite sur ces résumés déjà calculés, sans jamais rouvrir
l'encodeur. Chacune vous coûte quelques minutes, au lieu de recommencer un
entraînement complet à chaque fois.
\end{exemple}

Deuxième geste : on branche derrière lui une toute petite machine à apprendre,
qui, elle, a le droit d'apprendre. C'est elle qui aura besoin d'étiquettes —
mais très peu, puisqu'elle est minuscule.

\begin{definition}[La sonde][sonde]
Une \emph{sonde} est une petite machine à apprendre que l'on branche derrière
un encodeur gelé pour mesurer la qualité de ses représentations. Quand elle
se réduit à une seule \emph{couche} — c'est-à-dire une seule rangée de boutons
qui multiplient les nombres reçus et les additionnent — on parle de
\emph{sonde linéaire}. C'est un instrument de mesure, comme un thermomètre :
elle est volontairement faible, parce qu'une sonde faible ne peut pas
rattraper un mauvais résumé.
\end{definition}

\begin{schema}{Le principe de l'évaluation par sonde. L'encodeur porte un
cadenas : il est gelé, ses boutons ne bougent plus. Derrière lui, une toute
petite sonde apprend, elle, à répondre à une question précise. Si elle y
arrive avec très peu d'exemples, c'est que l'information était déjà bien
rangée dans la représentation.}[sonde]
  \begin{scope}[x=1cm, y=1cm]
    \draw[draw=mbSlate!55, fill=white, line width=0.7pt]
          (0,0.55) rectangle (1.1,1.55);
    \node[font=\sffamily\scriptsize, text=mbSlate] at (0.55,0.2) {une image};
    \draw[fleche] (1.25,1.05) -- (1.9,1.05);
    % encodeur gelé
    \node[bloc, minimum width=24mm, minimum height=13mm] (enc) at (3.2,1.05)
         {l'encodeur};
    % le cadenas, dessiné à la main : un corps et une anse
    \begin{scope}[shift={(3.2,2.15)}]
      \draw[draw=mbSlate, line width=0.9pt] (-0.16,0.1) arc (180:0:0.16);
      \draw[draw=mbSlate, fill=mbMist, line width=0.9pt]
            (-0.24,-0.28) rectangle (0.24,0.12);
      \fill[mbSlate] (0,-0.08) circle (0.05);
    \end{scope}
    \node[font=\sffamily\scriptsize\itshape, text=mbSlate] at (3.2,2.7)
         {gelé : plus aucun bouton ne bouge};
    \draw[fleche accent] (4.45,1.05) -- (5.15,1.05);
    \node[font=\sffamily\scriptsize, text=mbAmber!75!black]
         at (4.8,0.2) {la représentation};
    % la sonde, volontairement petite
    \node[bloc accent, minimum width=13mm, minimum height=8mm] (sonde)
         at (5.9,1.05) {\scriptsize la sonde};
    \draw[fleche] (6.6,1.05) -- (7.25,1.05);
    \node[font=\sffamily\small, text=mbInk, align=center] at (8.15,1.05)
         {la réponse\\ \scriptsize « chat »};
    \node[annotation, text width=3.1cm] at (8.0,2.35)
         {toute petite, et c'est\\ la seule qui apprend};
    \draw[trait annotation] (7.0,2.1) -- (6.15,1.55);
  \end{scope}
\end{schema}

\begin{intuition}
Pourquoi prendre une sonde volontairement faible ? Parce qu'on cherche à
mesurer la qualité du \emph{rangement}, pas la puissance de l'instrument de
mesure.

Imaginez deux ateliers. Dans le premier, les outils sont rangés par taille,
étiquetés, à portée de main. Dans le second, tout est jeté en vrac dans une
caisse. Envoyez un enfant chercher une clé de douze : dans le premier
atelier, il la trouve en dix secondes ; dans le second, jamais. Si vous
envoyez un mécanicien professionnel, il finira par la trouver dans les deux
cas — et vous n'aurez rien appris sur le rangement. La sonde, c'est l'enfant.
\end{intuition}

Nous reviendrons longuement sur ces mesures dans la partie~III, où nous
verrons comment on note vraiment une représentation. Pour l'instant, retenez
le principe : \emph{encodeur gelé, petite sonde, peu d'étiquettes}.

Ces deux gestes découpent le travail en deux temps. Le premier porte un nom
qu'il faut absolument connaître, car il revient dans tous les articles et
dans tout le reste du livre.

\begin{definition}[Le pré-entraînement][preentrainement]
Le \emph{pré-entraînement} est le long entraînement de départ, celui où la
machine joue à la tâche prétexte sur des montagnes de données brutes, sans
aucune étiquette. Le préfixe « pré » veut dire « avant » : c'est ce qui se
passe \emph{avant} qu'on lui demande quoi que ce soit d'utile. C'est de loin
la partie la plus coûteuse, et elle ne se fait qu'une seule fois.
\end{definition}

\begin{exemple}[Les deux temps, en chiffres]
Pour V-JEPA~2.1, le pré-entraînement, c'est cent trente-cinq mille pas sur
cent soixante-trois millions d'images et de vidéos, sans une seule étiquette.
C'est un travail que l'on ne fait qu'une fois, et les auteurs en ont publié
le résultat pour que personne n'ait à le recommencer.

L'utilisation, c'est tout autre chose : une petite sonde, quelques milliers
d'exemples étiquetés, et l'affaire est réglée. C'est ce deuxième temps, et
lui seul, que vous ferez tourner sur votre machine — y compris à la fin de ce
chapitre.
\end{exemple}

\begin{aretenir}
L'auto-supervision se déroule toujours en deux temps, et il ne faut jamais
les confondre :

\begin{enumerate}
  \item \textbf{Le pré-entraînement.} La machine joue au jeu des trous sur des
        montagnes de données brutes, sans la moindre étiquette. C'est long,
        c'est cher en calcul, et cela ne se fait qu'une fois. C'est ici que
        naît la représentation.
  \item \textbf{L'utilisation}, que l'on appelle aussi l'évaluation quand on
        cherche à noter le modèle. On gèle l'encodeur et l'on branche derrière
        une petite sonde, entraînée avec un tout petit nombre d'étiquettes,
        pour la tâche qui nous intéresse vraiment.
\end{enumerate}

Les étiquettes n'ont donc pas complètement disparu : elles ont été repoussées
à la toute fin, en quantité minuscule.
\end{aretenir}

Voyons maintenant à quelle échelle tout cela se joue dans le papier que nous
expliquons.

\section{À quelle échelle : les chiffres de V-JEPA 2.1}

Jusqu'ici, nous avons parlé de dix mille photos, puis de cent quarante-deux
millions. Mettons les vrais chiffres du papier sur la table, pour que vous
sachiez de quoi il retourne.

\begin{surlesdonnees}
Le jeu de données de V-JEPA~2.1 s'appelle \emph{VisionMix163M}. Son nom vient
tout simplement du nombre d'\emph{échantillons} qu'il contient — un
échantillon, c'est un élément de la collection, ici une image ou une vidéo :

\begin{center}
\footnotesize
\begin{tabular}{lrl}
\toprule
\textbf{Source} & \textbf{Échant.} & \textbf{Ce que c'est} \\
\midrule
LVD-142M    & $142{,}0$ M & des images choisies avec soin \\
YT-1B       & $19{,}0$ M  & des vidéos prises sur internet \\
HowTo100M   & $1{,}1$ M   & des vidéos qui expliquent \\
Kinetics    & $0{,}733$ M & des actions filmées de loin \\
SSv2        & $0{,}168$ M & des gestes filmés de près \\
\midrule
\textbf{Total} & $\approx 163$ M & d'où le nom \\
\bottomrule
\end{tabular}
\end{center}

Faites l'addition vous-même :
$142 + 19 + 1{,}1 + 0{,}733 + 0{,}168 = 163{,}0$ millions. Mises bout à bout,
ces vidéos durent environ $1{,}7$ million d'heures, dont $1\,600\,000$ pour la
seule source YT-1B.

Et voici le point qui doit rester : \textbf{aucun de ces cent soixante-trois
millions d'échantillons n'a été étiqueté à la main pour cet entraînement}.
Le calcul de la section~2 nous a montré qu'il aurait fallu, au minimum, cent
vingt-huit ans pour les images seules. Ce livre entier raconte comment on
s'en passe.
\end{surlesdonnees}

\begin{remarque}
Vous avez peut-être remarqué que certains de ces jeux de données possèdent
bel et bien des étiquettes : Kinetics et SSv2 sont même nés pour cela. Elles
existent, mais V-JEPA~2.1 ne les regarde pas pendant le pré-entraînement. Il
ne prend que les pixels. Les étiquettes ne réapparaissent qu'à la toute fin,
pour entraîner les sondes et pour noter les résultats.
\end{remarque}

Assez parlé. Il est temps de faire tourner tout cela pour de vrai, sur votre
machine, avec un programme que vous pouvez lire en entier.

\section{Passons au code}

Nous allons construire le plus petit exemple d'auto-supervision possible, et
vérifier de nos yeux qu'il fonctionne. Le programme tient en une page, et il
suit exactement les deux temps de l'encadré précédent : d'abord jouer au jeu
des trous sans aucune étiquette, ensuite geler et brancher une sonde.

Voici notre monde miniature. Chaque image fait $8 \times 8$ cases, et contient
une seule barre : soit une barre \emph{couchée}, c'est-à-dire horizontale,
soit une barre \emph{debout}, c'est-à-dire verticale. Elle est posée
n'importe où. Nous dirons toujours « couchée » et « debout », y compris dans
le programme.

\begin{schema}{Notre monde miniature et le jeu que nous allons y jouer.
Chaque image contient une barre, couchée ou debout. On coupe l'image en deux :
la moitié gauche est donnée à la machine, la moitié droite est cachée et
devient la bonne réponse. La machine ne saura jamais qu'il existe des mots
comme « couchée » et « debout ».}[monde-barres]
  \begin{scope}[x=1cm, y=1cm]
    % --- image 1 : barre horizontale ---
    \begin{scope}[shift={(0,0)}]
      \foreach \mbi in {0,...,7}{
        \foreach \mbj in {0,...,7}{
          \draw[draw=mbSlate!30, fill=white, line width=0.3pt]
                (\mbi*0.3,\mbj*0.3) rectangle (\mbi*0.3+0.3,\mbj*0.3+0.3);
        }
      }
      \foreach \mbi in {0,...,7}{
        \draw[draw=mbSlate!30, fill=mbBlue!45, line width=0.3pt]
              (\mbi*0.3,1.5) rectangle (\mbi*0.3+0.3,1.8);
      }
      \node[font=\sffamily\scriptsize, text=mbSlate] at (1.2,-0.4)
           {barre couchée};
    \end{scope}
    % --- image 2 : barre verticale ---
    \begin{scope}[shift={(3.1,0)}]
      \foreach \mbi in {0,...,7}{
        \foreach \mbj in {0,...,7}{
          \draw[draw=mbSlate!30, fill=white, line width=0.3pt]
                (\mbi*0.3,\mbj*0.3) rectangle (\mbi*0.3+0.3,\mbj*0.3+0.3);
        }
      }
      \foreach \mbj in {0,...,7}{
        \draw[draw=mbSlate!30, fill=mbTeal!45, line width=0.3pt]
              (0.6,\mbj*0.3) rectangle (0.9,\mbj*0.3+0.3);
      }
      \node[font=\sffamily\scriptsize, text=mbSlate] at (1.2,-0.4)
           {barre debout};
    \end{scope}
    % --- la coupe en deux ---
    \begin{scope}[shift={(6.6,0)}]
      \foreach \mbi in {0,...,7}{
        \foreach \mbj in {0,...,7}{
          \draw[draw=mbSlate!30, fill=white, line width=0.3pt]
                (\mbi*0.3,\mbj*0.3) rectangle (\mbi*0.3+0.3,\mbj*0.3+0.3);
        }
      }
      \foreach \mbi in {0,...,7}{
        \draw[draw=mbSlate!30, fill=mbBlue!45, line width=0.3pt]
              (\mbi*0.3,0.9) rectangle (\mbi*0.3+0.3,1.2);
      }
      \fill[mbSlate!45, opacity=0.85] (1.2,0) rectangle (2.4,2.4);
      \draw[draw=mbSlate, line width=1pt, dashed] (1.2,-0.1) -- (1.2,2.5);
      \node[font=\sffamily\scriptsize, text=mbSlate] at (0.6,-0.4)
           {donné};
      \node[font=\sffamily\scriptsize, text=white] at (1.8,1.2) {caché};
      \node[font=\sffamily\scriptsize, text=mbSlate] at (1.8,-0.4)
           {à deviner};
    \end{scope}
    \draw[accolade] (6.6,2.6) -- (9.0,2.6);
    \node[annotation, text width=3.4cm] at (7.8,3.4)
         {la question et la réponse\\ sortent de la même image :\\
          aucune étiquette};
  \end{scope}
\end{schema}

Commençons par fabriquer ces images. Le programme ci-dessous ne contient rien
que nous n'ayons expliqué : des tenseurs, comme au
chapitre~\ref{chap:image-video}, et des tirages au hasard.

\begin{code}{python}[Fabriquer le monde miniature]
import torch

torch.manual_seed(0)          # pour que vous obteniez les memes nombres

def fabrique_images(n):
    """n images de 8x8 : soit une barre couchee, soit une barre debout."""
    images = torch.zeros(n, 8, 8)
    genre = torch.randint(0, 2, (n,))    # 0 = couchee, 1 = debout
    place = torch.randint(0, 8, (n,))    # a quelle ligne, ou a quelle colonne
    for i in range(n):
        if genre[i] == 0:
            images[i, place[i], :] = 1.0     # une ligne entiere
        else:
            images[i, :, place[i]] = 1.0     # une colonne entiere
    return images, genre.float()

images, genre = fabrique_images(2000)
gauche = images[:, :, :4].reshape(-1, 32)   # 8 lignes x 4 colonnes = 32
droite = images[:, :, 4:].reshape(-1, 32)   # l'autre moitie
\end{code}

Deux mots sur les deux dernières lignes, car elles contiennent un geste que
nous referons souvent. L'écriture \co{python}{images[:, :, :4]} veut dire :
« garde toutes les images, toutes les lignes, mais seulement les quatre
premières colonnes » — c'est notre moitié gauche. Quant à
\co{python}{reshape(-1, 32)}, il met chaque petite grille de $8 \times 4$ à
plat, en une seule rangée de $8 \times 4 = 32$ nombres, comme si l'on
découpait les lignes d'un tableau pour les recoller bout à bout. Le $-1$ dit
simplement à la machine : « débrouille-toi pour le nombre d'images, tu sais
compter ».

Notez bien que \co{python}{genre} existe. C'est une \emph{variable},
c'est-à-dire un nom que nous donnons à une valeur pour pouvoir la retrouver
plus loin dans le programme ; celle-ci retient, pour chaque image, si la
barre était couchée ou debout. Nous allons pourtant l'ignorer complètement
pendant tout le pré-entraînement : elle attendra sagement la toute fin, pour
la sonde. C'est notre façon d'imiter la réalité : les étiquettes existent
quelque part dans le monde, mais on ne s'en sert presque pas.

Construisons maintenant les deux pièces du modèle. La première, nous la
connaissons : c'est l'encodeur, qui résume les trente-deux nombres de la
moitié gauche en seize nombres. La seconde est nouvelle, et il faut lui
donner son nom.

\begin{definition}[Le prédicteur][predicteur]
Le \emph{prédicteur} est la pièce qui part du résumé fabriqué par l'encodeur
et qui tente de deviner ce qui a été caché. C'est lui qui joue, à proprement
parler, au jeu des trous. À la fin de l'entraînement, c'est pourtant
l'encodeur que l'on garde précieusement ; le prédicteur, lui, n'a servi qu'à
faire travailler son voisin. Nous verrons dans la partie~V qu'il sait tout de
même rendre un dernier service, quand on veut deviner la suite d'une vidéo.
\end{definition}

Chez nous, ce prédicteur part des seize nombres du résumé et tente de
reconstituer les trente-deux nombres cachés.

Trois outils de PyTorch apparaissent pour la première fois. Prenons-les un par
un, car nous les retrouverons dans tout le livre.

\begin{enclair}
\co{python}{torch.nn.Linear(32, 16)} fabrique une \emph{couche} : une rangée
de boutons qui prend trente-deux nombres en entrée, les multiplie chacun par
un poids, additionne le tout, et rend seize nombres. C'est exactement la
machine à bonbons du chapitre~\ref{chap:apprendre}, en beaucoup plus large :
au lieu des deux boutons $a$ et $b$, celle-ci en possède
$32 \times 16 + 16 = 528$.

\co{python}{torch.nn.ReLU()} est une règle minuscule appliquée après la
couche : tout nombre négatif est remplacé par zéro, les autres passent tels
quels. Sans elle, empiler deux couches ne servirait à rien, car deux
multiplications à la suite se ramènent toujours à une seule. C'est cette
petite règle qui permet à un modèle d'apprendre autre chose que des droites.
Quant à \mbox{\co{python}{Sequential}}, il sert simplement à mettre ces deux
pièces bout à bout.

\co{python}{torch.optim.Adam} remplace la ligne
\co{python}{a -= pas * a.grad} que nous écrivions à la main. C'est la même
descente de gradient, avec une amélioration : Adam règle tout seul la
longueur du pas de chaque bouton, en se souvenant des pentes passées. Le
réglage \co{python}{lr=0.01} est le taux d'apprentissage de départ, celui du
chapitre~\ref{chap:apprendre}. Restent deux lignes de service.
\co{python}{optim.step()} bouge tous les boutons d'un coup, et
\co{python}{optim.zero_grad()} efface les pentes avant le tour suivant. Ce
sont les deux gestes que nous écrivions bouton par bouton au chapitre
précédent.
\end{enclair}

Ces trois outils paraissent abstraits tant qu'on ne les a pas vus tourner sur
de vrais nombres. Faisons-le en tout petit.

\begin{exemple}[Une couche minuscule, calculée à la main]
Prenons une couche qui reçoit deux nombres et n'en rend qu'un seul. Elle
possède donc deux poids — disons $0{,}5$ et $-2$ — plus un décalage de $4$,
qui est exactement le bouton $b$ de la machine à bonbons. Donnons-lui les
nombres $3$ et $1$ :
\[
  0{,}5 \times 3 \;+\; (-2) \times 1 \;+\; 4
  \;=\; 1{,}5 - 2 + 4 \;=\; 3{,}5 .
\]
La couche rend donc $3{,}5$, et la règle ReLU le laisse passer tel quel
puisqu'il est positif.

Refaisons le calcul avec les nombres $1$ et $4$ :
\[
  0{,}5 \times 1 \;+\; (-2) \times 4 \;+\; 4
  \;=\; 0{,}5 - 8 + 4 \;=\; -3{,}5 .
\]
Cette fois le résultat est négatif : ReLU le remplace par $0$, et c'est ce
zéro qui poursuivra le voyage. Voilà tout le mystère. Notre vrai encodeur
fait la même chose, mais avec trente-deux nombres à l'entrée, seize à la
sortie, et $528$ boutons au lieu de trois.
\end{exemple}

\begin{code}{python}[Le pré-entraînement, sans la moindre étiquette]
encodeur = torch.nn.Sequential(torch.nn.Linear(32, 16), torch.nn.ReLU())
predicteur = torch.nn.Linear(16, 32)

params = list(encodeur.parameters()) + list(predicteur.parameters())
optim = torch.optim.Adam(params, lr=0.01)

for etape in range(1, 601):
    erreur = ((predicteur(encodeur(gauche)) - droite) ** 2).mean()
    optim.zero_grad()
    erreur.backward()
    optim.step()
    if etape in (1, 100, 300, 600):
        print(f"etape {etape:3d} : erreur = {erreur.item():.4f}")
\end{code}

Vous reconnaissez la boucle en cinq gestes du chapitre~\ref{chap:apprendre},
à un détail près : la bonne réponse, \co{python}{droite}, n'a été écrite par
personne. Elle a été découpée dans l'image elle-même.

Voici ce que le programme affiche :

\begin{sortie}
etape   1 : erreur = 0.1648
etape 100 : erreur = 0.0487
etape 300 : erreur = 0.0483
etape 600 : erreur = 0.0483
\end{sortie}

L'erreur tombe de $0{,}1648$ à $0{,}0483$, puis ne bouge plus. Cette
stagnation n'est pas une panne, et c'est même le passage le plus instructif
du chapitre.

\begin{piege}
Quand une barre est debout dans la moitié \emph{droite} de l'image, la moitié
gauche est entièrement vide. On demande alors à la machine de deviner une
barre dont absolument rien, dans ce qu'elle voit, n'indique la position.
Vous-même, devant une moitié gauche toute blanche, ne feriez pas mieux.

Dans notre tirage, cela arrive pour $25{,}8\,\%$ des images. Il existe donc
un plancher sous lequel l'erreur ne peut pas descendre, et le modèle l'a
simplement atteint. Ne confondez jamais « le modèle a cessé de progresser »
avec « le modèle est mauvais » : il faut d'abord se demander si ce qu'on lui
demande est seulement possible.
\end{piege}

Passons au deuxième temps : geler, brancher une sonde, et compter les points.
Nous n'accorderons à la sonde que \emph{vingt} exemples étiquetés — c'est
volontairement très peu.

Une dernière nouveauté, et ce sera la seule. Jusqu'ici, notre erreur était la
moyenne des carrés des écarts, parce que nous prédisions des nombres. Ici, la
sonde doit répondre à une question par oui ou par non : cette barre est-elle
debout ? Pour ce genre de question, on emploie une autre mesure d'erreur,
\co{python}{binary_cross_entropy_with_logits}. Son nom fait peur, mais son
idée tient en une phrase : elle punit d'autant plus la machine qu'elle s'est
trompée \emph{avec assurance}. Répondre « oui » à moitié convaincu quand la
réponse était non coûte peu ; répondre « oui, absolument certain » quand la
réponse était non coûte très cher.

La sonde rend un seul nombre par image. On lit ce nombre ainsi : s'il est
plus grand que zéro, la machine dit « debout » ; sinon, elle dit « couchée ».
C'est le rôle du test \co{python}{> 0} à l'avant-dernière ligne. Au passage,
le mot \emph{logits}, dans le nom de la mesure d'erreur, désigne simplement
ce nombre brut, celui qu'on lit avant de le traduire en « oui » ou en
« non ». Quant au petit \co{python}{squeeze(1)}, il ne fait que retirer un
axe devenu inutile, pour que les formes des deux tenseurs se correspondent :
rien de plus qu'un rangement, au sens du chapitre~\ref{chap:image-video}.

\begin{code}{python}[Geler, puis brancher une sonde sur vingt étiquettes]
def sonde(entrees, cibles, entrees_test, cibles_test):
    """Une seule couche, qui apprend a repondre : couchee ou debout ?"""
    tete = torch.nn.Linear(entrees.shape[1], 1)
    opt = torch.optim.Adam(tete.parameters(), lr=0.05)
    for _ in range(300):
        p = tete(entrees).squeeze(1)
        perte = torch.nn.functional.binary_cross_entropy_with_logits(p, cibles)
        opt.zero_grad()
        perte.backward()
        opt.step()
    with torch.no_grad():
        prediction = (tete(entrees_test).squeeze(1) > 0).float()
        return (prediction == cibles_test).float().mean().item()

test_images, test_genre = fabrique_images(1000)
test_gauche = test_images[:, :, :4].reshape(-1, 32)

with torch.no_grad():                  # l'encodeur est gele : aucun gradient
    h = encodeur(gauche)
    h_test = encodeur(test_gauche)

torch.manual_seed(1)
print("sonde sur la representation :",
      f"{sonde(h[:20], genre[:20], h_test, test_genre)*100:.1f} %")
torch.manual_seed(1)
print("sonde sur les pixels bruts  :",
      f"{sonde(gauche[:20], genre[:20], test_gauche, test_genre)*100:.1f} %")
\end{code}

Et voici le résultat, mesuré sur mille images que la machine n'avait jamais
vues :

\begin{sortie}[Le verdict]
sonde sur la representation : 80.4 %
sonde sur les pixels bruts  : 61.8 %
\end{sortie}

Arrêtons-nous, car ces deux lignes résument tout le chapitre.

La même sonde, avec les mêmes vingt étiquettes et le même nombre d'étapes,
obtient $80{,}4\,\%$ de bonnes réponses quand on la branche sur la
représentation, contre $61{,}8\,\%$ quand on la branche directement sur les
pixels. La différence ne vient pas de la sonde, qui est identique dans les
deux cas. Elle vient de ce que l'encodeur a appris \emph{en jouant à un jeu
qui n'avait rien à voir}. Personne ne lui a jamais montré le mot « debout »
ni le mot « couchée ».

\begin{remarque}
Et $80{,}4\,\%$, est-ce beaucoup ? Refaisons le calcul du piège précédent. Dans
$25{,}8\,\%$ des cas, la moitié gauche est vide et personne ne peut répondre
autrement qu'en tirant à pile ou face, ce qui rapporte la moitié des points,
soit environ $12{,}9\,\%$. Le meilleur score atteignable est donc de l'ordre de
$74{,}2 + 12{,}9 \approx 87\,\%$. Avec $80{,}4\,\%$, notre encodeur n'est pas
loin du maximum possible, et il l'a obtenu avec vingt étiquettes.
\end{remarque}

Vous n'avez pas à recopier ces morceaux à la main : le programme entier, tel
que nous venons de le construire, est rangé dans le fichier
\term{scripts/03-auto-supervision.py} qui accompagne ce livre. Lancez-le par
la commande \co{text}{python 03-auto-supervision.py} et vous devez retrouver,
à la décimale près, les nombres imprimés ci-dessus.

Si vous obtenez autre chose, regardez d'abord la toute première ligne,
\co{python}{torch.manual_seed(0)}. Un ordinateur ne sait pas tirer au sort
comme on lance un dé : il fabrique des nombres qui \emph{ressemblent} à du
hasard, en partant d'un nombre de départ qu'on appelle la \emph{graine}.
Avec la même graine, il refait exactement les mêmes tirages, dans le même
ordre. C'est ce qui vous permet de retrouver nos barres, nos vingt étiquettes
et nos résultats à l'identique ; changez le $0$ en $1$, et tout le monde
miniature change avec.

\begin{enia}
Changez maintenant les échelles, et vous avez V-JEPA~2.1. Les images de
$8 \times 8$ deviennent des vidéos de soixante-quatre images en
$384 \times 384$. Les seize nombres de la représentation deviennent des
milliers. Les six cents étapes deviennent cent trente-cinq mille. Les deux
mille images deviennent cent soixante-trois millions d'images et de vidéos.
Le prédicteur, lui, devient une machine à calculer bien plus grosse, que nous
démonterons pièce par pièce le moment venu.

Mais la structure du programme, elle, ne change pas d'une ligne : cacher,
deviner, mesurer, corriger — puis geler et brancher une sonde.
\end{enia}

\begin{aretenir}
Les trois idées du chapitre, en trois phrases.

\begin{enumerate}
  \item Les étiquettes coûtent trop cher pour être produites à l'échelle des
        vidéos : il fallait une autre source de bonnes réponses.
  \item Cette source, c'est la donnée elle-même : on cache un morceau, et le
        morceau caché fait office de bonne réponse.
  \item Le jeu ne compte pas ; ce qui compte, c'est la représentation que la
        machine fabrique en y jouant, et que l'on mesure avec une petite
        sonde derrière un encodeur gelé.
\end{enumerate}
\end{aretenir}

Il reste pourtant une question que nous avons soigneusement contournée, et
elle est bien plus profonde qu'elle n'en a l'air. Dans notre petit programme,
nous avons demandé à la machine de deviner les \emph{pixels} cachés : les
trente-deux nombres de la moitié droite, un par un. Est-ce vraiment ce qu'il
faut demander ?

Réfléchissez un instant à ce que cela implique. Pour redessiner exactement
les pixels d'une zone cachée dans une vraie vidéo, il faudrait deviner chaque
brin d'herbe, chaque grain de la peau, chaque reflet dans une flaque. Or
personne au monde ne peut prédire cela, et surtout, personne n'en a besoin.
La machine passerait le plus clair de son temps sur un travail inutile.

Il existe une autre possibilité, beaucoup plus fine : lui demander de deviner
non pas les pixels cachés, mais \emph{le résumé} de la zone cachée. Non pas
l'herbe brin par brin, mais l'idée « ici, il y a de l'herbe ». C'est-à-dire
prédire dans l'espace latent que nous venons de définir.

Cette question — \emph{dans quel espace faut-il deviner : celui des pixels,
ou celui des idées ?} — est le cœur de tout le papier que ce livre
explique. Sa réponse porte un nom, JEPA, et c'est le sujet de la partie~II.

Avant d'y arriver, il nous manque encore une pièce. Nous savons ce que la
machine doit apprendre, mais nous n'avons pas encore vu \emph{avec quoi} elle
regarde. Quel est l'œil de la machine, celui qui découpe une vidéo en
morceaux et les compare entre eux ? C'est le sujet du chapitre suivant.

\phantomsection
\section*{Exercices}
\addcontentsline{toc}{section}{Exercices}

\begin{exercice}[Le prix d'une étiquette]
Vous devez faire étiqueter $500\,000$ images. Chaque étiquette prend huit
secondes. Vous employez cinq personnes, qui travaillent chacune sept heures
par jour. Combien de jours faudra-t-il ?
\end{exercice}

\begin{solution}
Le travail total vaut $500\,000 \times 8 = 4\,000\,000$ secondes, soit
$4\,000\,000 / 3\,600 \approx 1\,111$ heures. Chaque jour, les cinq personnes
fournissent ensemble $5 \times 7 = 35$ heures. Il faut donc

\[
  \frac{1\,111}{35} \approx 32 \text{ jours.}
\]

Un mois de travail à cinq, pour un demi-million d'images seulement. Multipliez
par deux cent quatre-vingt-quatre pour atteindre les cent quarante-deux
millions d'images de V-JEPA~2.1, et vous retrouvez l'ordre de grandeur de la
section~2.
\end{solution}

\begin{exercice}[Fabriquer une tâche prétexte]
On vous donne un million de vidéos sans aucune étiquette. Proposez deux jeux
que la machine pourrait jouer toute seule, et dites à chaque fois d'où sort
la bonne réponse.
\end{exercice}

\begin{solution}
Premier jeu : remettre les images dans l'ordre. On mélange l'ordre des images
d'un clip et l'on demande de le rétablir ; la bonne réponse est l'ordre
d'origine, que l'on connaît puisque c'est nous qui l'avons mélangé.

Deuxième jeu : deviner la suite. On coupe le clip en deux et l'on demande de
prédire la seconde moitié à partir de la première ; la bonne réponse est la
seconde moitié, que l'on possède déjà.

Dans les deux cas, la règle est la même : on abîme volontairement la donnée
d'une façon dont on garde la trace, et cette trace est la bonne réponse.
\end{solution}

\begin{exercice}[Le taux de compression]
Un clip de $16$ images en $256 \times 256$ pixels couleur est résumé par un
encodeur en $512$ nombres. Par combien l'information a-t-elle été divisée ?
\end{exercice}

\begin{solution}
Le clip contient $16 \times 256 \times 256 \times 3 = 3\,145\,728$ nombres,
comme nous l'avions calculé au chapitre~\ref{chap:image-video}. Le rapport
vaut

\[
  \frac{3\,145\,728}{512} = 6\,144 .
\]

Le résumé est six mille cent quarante-quatre fois plus petit que le clip. Toute la
question est de savoir si l'on a jeté les bons nombres — et c'est précisément
ce que mesure une sonde.
\end{solution}

\begin{exercice}[Lire un résultat de sonde]
Deux encodeurs sont évalués avec exactement la même sonde et les mêmes
étiquettes. Le premier obtient $78\,\%$, le second $62\,\%$. Que peut-on en
conclure, et que ne peut-on \emph{pas} en conclure ?
\end{exercice}

\begin{solution}
On peut conclure que le premier encodeur range mieux, dans ses
représentations, l'information dont cette tâche précise a besoin : une simple
sonde y accède plus facilement.

On ne peut pas conclure que le premier encodeur est meilleur en général. Il
pourrait être moins bon sur une autre tâche, par exemple mesurer des
distances plutôt que reconnaître des objets. Nous verrons dans la partie~V
que c'est exactement ce qui arrive à V-JEPA~2 : excellent pour reconnaître une
action, faible pour décrire chaque pixel. Un chiffre de sonde ne vaut que
pour la question posée.
\end{solution}

\begin{exercice}[Le plancher d'erreur]
Dans le programme de ce chapitre, l'erreur cesse de descendre à $0{,}0483$.
Un collègue vous dit : « ton modèle est bloqué, augmente le taux
d'apprentissage ». Que lui répondez-vous ?
\end{exercice}

\begin{solution}
Qu'il faut d'abord vérifier si la tâche est possible. Ici, dans $25{,}8\,\%$
des cas, la moitié gauche de l'image est entièrement vide alors qu'une barre
se cache à droite : aucune information ne permet de la placer. Une partie de
l'erreur est donc irréductible, et le modèle a simplement atteint ce plancher.

Augmenter le taux d'apprentissage ne servirait à rien, et risquerait même de
faire exploser l'entraînement, comme nous l'avons vu au
chapitre~\ref{chap:apprendre}. La bonne réaction est de se demander ce que la
machine \emph{pourrait} savoir, avant de lui reprocher de ne pas le savoir.
\end{solution}
